% ============================================================
\setpartpreamble{\partintro
  {How LINQ becomes SQL and how a plugin translates its own functions and members.}
  {
    \item Chapter~\ref{ch:the-query-pipeline}: The query pipeline
    \item Chapter~\ref{ch:translating-methods-and-members}: Translating methods and members
    \item Chapter~\ref{ch:preventing-client-side-evaluation}: Preventing client-side
    evaluation
    \item Chapter~\ref{ch:further-query-extension-points}: Further query extension points
  }}
\part{Queries}\label{part:IV}

% ------------------------------------------------------------
\chapter{The query pipeline}\label{ch:the-query-pipeline}

A LINQ query passes through several stages in EF Core before SQL is produced. Knowing where
each stage sits tells you where a plugin has to hook in:

\begin{flowdiagram}
  \flowstep{Funcletizer}{Evaluates sub-expressions without data reference in .NET and turns
  them into parameters. Plugin: IEvaluatableExpressionFilterPlugin}
  \flowstep{Preprocessing}{Expands navigations, inserts query filters, normalizes. Plugin:
  IQueryExpressionInterceptor}
  \flowstep{Queryable translation}{Where, Select, OrderBy etc. become a SelectExpression}
  \flowstep{SQL translation}{Expressions become SqlExpressions. Plugins: method, member and
  aggregate translators}
  \flowstep{Postprocessing}{Nullability optimization, parameter-dependent adjustments}
  \flowstep{SQL generation}{QuerySqlGenerator writes the SQL text}
  \flowstep{Materialization}{Rows become objects. Plugin: IMaterializationInterceptor}
\end{flowdiagram}

The compiled query is cached. Translators therefore run only once per query shape, not on
every execution. Values that change per execution (parameters, session settings) must
therefore appear in SQL as parameters or as a database function, never as a constant
inserted at translation time. That is exactly why the sample translates
\code{Current\allowbreak{}Tenant(\allowbreak{})} into a call to
\code{current\_\allowbreak{}setting(\allowbreak{}.\allowbreak{}.\allowbreak{}.\allowbreak{})}
instead of inserting the current tenant.

% ------------------------------------------------------------
\chapter{Translating methods and members}\label{ch:translating-methods-and-members}

\section{The EF.Functions pattern}

Database functions without a .NET counterpart are declared as extension methods on
\code{Db\allowbreak{}Functions}. They are never executed; the body only throws an exception
in case someone calls them outside of LINQ. The declaration lives in the core package, the
translation in the provider package:

\begin{sourcecode}[Tenancy\allowbreak{}Db\allowbreak{}Functions\allowbreak{}Extensions.\allowbreak{}cs (\allowbreak{}core package)]{csharp}
namespace Microsoft.EntityFrameworkCore;

public static class TenancyDbFunctionsExtensions
{
    public static TenantId CurrentTenant(this DbFunctions _)
        => throw new InvalidOperationException(
            $"{nameof(CurrentTenant)} can only be used in LINQ to Entities queries.");
}
\end{sourcecode}

\section{A method call translator}

An \code{IMethod\allowbreak{}Call\allowbreak{}Translator} receives every method call EF does
not know itself. It checks whether it is responsible and then builds a SQL expression tree
using \code{ISql\allowbreak{}Expression\allowbreak{}Factory}:

\begin{sourcecode}[Query/\allowbreak{}Current\allowbreak{}Tenant\allowbreak{}Translator.\allowbreak{}cs]{csharp}
public class CurrentTenantTranslator : IMethodCallTranslator
{
    private static readonly MethodInfo CurrentTenantMethod =
        typeof(TenancyDbFunctionsExtensions).GetRuntimeMethod(
            nameof(TenancyDbFunctionsExtensions.CurrentTenant), [typeof(DbFunctions)])!;

    private readonly ISqlExpressionFactory _sql;
    private readonly IRelationalTypeMappingSource _typeMappingSource;
    private readonly ITenancySingletonOptions _options;

    public CurrentTenantTranslator(
        ISqlExpressionFactory sql,
        IRelationalTypeMappingSource typeMappingSource,
        ITenancySingletonOptions options)
    {
        // ... ArgumentNullException.ThrowIfNull for all parameters ...
        _sql = sql;
        _typeMappingSource = typeMappingSource;
        _options = options;
    }

    public virtual SqlExpression? Translate(
        SqlExpression? instance, MethodInfo method,
        IReadOnlyList<SqlExpression> arguments,
        IDiagnosticsLogger<DbLoggerCategory.Query> logger)
    {
        if (method != CurrentTenantMethod)
            return null;                          // not responsible -> next translator

        var text = _typeMappingSource.FindMapping(typeof(string))!;
        var boolean = _typeMappingSource.FindMapping(typeof(bool))!;

        var currentSetting = _sql.Function(
            "current_setting",
            [
                _sql.Constant(_options.Settings.SettingName, text),
                _sql.Constant(true, boolean),
            ],
            nullable: true,
            argumentsPropagateNullability: [true, false],
            typeof(string), text);

        var nullIfEmpty = _sql.Function(
            "nullif",
            [currentSetting, _sql.Constant(string.Empty, text)],
            nullable: true,
            argumentsPropagateNullability: [true, false],
            typeof(string), text);

        return _sql.Convert(nullIfEmpty, typeof(TenantId), TenantIdTypeMapping.Default);
    }
}
\end{sourcecode}

This turns \code{o.\allowbreak{}Tenant\allowbreak{}Id ==
EF.\allowbreak{}Functions.\allowbreak{}Current\allowbreak{}Tenant(\allowbreak{})} into:

\begin{sourcecode}[Generated SQL (\allowbreak{}excerpt)]{sql}
WHERE o."TenantId" = CAST(nullif(current_setting('app.tenant_id', TRUE), '') AS uuid)
\end{sourcecode}

\begin{datatable}{H{0.60}H{1.40}}
  \textbf{Detail} & \textbf{Meaning} \\
  \tablehead
  \code{return null} & Contract for “not responsible”. Never throw for foreign methods; the
  next translator gets its turn. \\
  \tablerule
  \code{arguments} & For extension methods, \code{arguments[0]} is the \code{this} parameter
  (here \code{Db\allowbreak{}Functions}). \\
  \tablerule
  \code{nullable} & Can the result be NULL? EF uses this to optimize \code{IS NULL} checks.
  Wrong values lead to wrong results. \\
  \tablerule
  \code{arguments\allowbreak{}Propagate\allowbreak{}Nullability} & Is the result NULL when
  this argument is NULL? For \code{nullif(\allowbreak{}x,\allowbreak{} '')}: yes for x, no
  for the empty string. \\
  \tablerule
  Type mappings & Every SQL expression needs a mapping. It determines literals, parameters
  and casts. \code{Convert} with the TenantId mapping produces
  \code{CAST(\allowbreak{}.\allowbreak{}.\allowbreak{}.\allowbreak{} AS uuid)}. \\
  \tablerule
  \code{Constant(\allowbreak{}.\allowbreak{}.\allowbreak{}.\allowbreak{})} & Constants
  become literals via the mapping and are escaped correctly, including the setting name. \\
\end{datatable}

\section{A member translator}

For properties and fields there is \code{IMember\allowbreak{}Translator}. The sample
translates \code{tenant\allowbreak{}Id.\allowbreak{}Is\allowbreak{}Empty}. The constant gets
the column's type mapping so that the literal comes out in the right format:

\begin{sourcecode}[Query/\allowbreak{}Tenant\allowbreak{}Id\allowbreak{}Member\allowbreak{}Translator.\allowbreak{}cs]{csharp}
public class TenantIdMemberTranslator : IMemberTranslator
{
    private static readonly PropertyInfo IsEmptyProperty =
        typeof(TenantId).GetProperty(nameof(TenantId.IsEmpty))!;

    private readonly ISqlExpressionFactory _sql;

    public TenantIdMemberTranslator(ISqlExpressionFactory sql)
    {
        ArgumentNullException.ThrowIfNull(sql);

        _sql = sql;
    }

    public virtual SqlExpression? Translate(
        SqlExpression? instance, MemberInfo member, Type returnType,
        IDiagnosticsLogger<DbLoggerCategory.Query> logger)
    {
        if (instance is null || !Equals(member, IsEmptyProperty))
            return null;

        return _sql.Equal(instance, _sql.Constant(TenantId.Empty, instance.TypeMapping));
    }
}
\end{sourcecode}

\section{Translator plugins}

Translators are not registered directly, but through plugin classes that return a list. The
plugins are \textbf{scoped} (appendix~\ref{app:A}) and can therefore inject scoped services
such as \code{ISql\allowbreak{}Expression\allowbreak{}Factory}:

\begin{sourcecode}[Query/\allowbreak{}Tenancy\allowbreak{}Translator\allowbreak{}Plugins.\allowbreak{}cs]{csharp}
public class TenancyMethodCallTranslatorPlugin : IMethodCallTranslatorPlugin
{
    public TenancyMethodCallTranslatorPlugin(
        ISqlExpressionFactory sql,
        IRelationalTypeMappingSource typeMappingSource,
        ITenancySingletonOptions options)
    {
        Translators = [new CurrentTenantTranslator(sql, typeMappingSource, options)];
    }

    public virtual IEnumerable<IMethodCallTranslator> Translators { get; }
}

public class TenancyMemberTranslatorPlugin : IMemberTranslatorPlugin
{
    public TenancyMemberTranslatorPlugin(ISqlExpressionFactory sql)
    {
        Translators = [new TenantIdMemberTranslator(sql)];
    }

    public virtual IEnumerable<IMemberTranslator> Translators { get; }
}
\end{sourcecode}

\begin{tip}
  \code{query.\allowbreak{}To\allowbreak{}Query\allowbreak{}String(\allowbreak{})} produces
  a query's SQL without executing it and without a database connection. It is the most
  important tool for developing and testing translators (chapter~\ref{ch:test-strategy}).
\end{tip}

% ------------------------------------------------------------
\chapter{Preventing client-side evaluation}\label{ch:preventing-client-side-evaluation}

The first stage of the pipeline, the funcletizer, evaluates every sub-expression in .NET up
front that has no reference to query data. \code{o.\allowbreak{}Total > limit * 2} thus
becomes a parameter. For
\code{EF.\allowbreak{}Functions.\allowbreak{}Current\allowbreak{}Tenant(\allowbreak{})} this
would be fatal: the expression has no data reference, so EF would execute the method, which
throws its exception.

EF knows only a few exceptions. In the source of
\code{Evaluatable\allowbreak{}Expression\allowbreak{}Filter} (EF Core 10) they are
\code{Date\allowbreak{}Time.\allowbreak{}Now},
\code{Date\allowbreak{}Time.\allowbreak{}Utc\allowbreak{}Now},
\code{Date\allowbreak{}Time.\allowbreak{}Today}, their
\code{Date\allowbreak{}Time\allowbreak{}Offset} counterparts,
\code{Guid.\allowbreak{}New\allowbreak{}Guid(\allowbreak{})},
\code{Random.\allowbreak{}Next(\allowbreak{}.\allowbreak{}.\allowbreak{}.\allowbreak{})} and
all methods declared in EF's own class
\code{Db\allowbreak{}Functions\allowbreak{}Extensions}. Custom EF.Functions methods live in
a different class and are \textbf{not} covered. After that, EF asks all registered plugins:

\begin{sourcecode}[Core package: Tenancy\allowbreak{}Db\allowbreak{}Functions\allowbreak{}Attribute.\allowbreak{}cs,\allowbreak{} Query/\allowbreak{}Tenancy\allowbreak{}Evaluatable…Plugin.\allowbreak{}cs]{csharp}
[AttributeUsage(AttributeTargets.Class, Inherited = false)]
public sealed class TenancyDbFunctionsAttribute : Attribute;

public class TenancyEvaluatableExpressionFilterPlugin : IEvaluatableExpressionFilterPlugin
{
    public virtual bool IsEvaluatableExpression(Expression expression)
        => expression is not MethodCallExpression { Method.DeclaringType: { } type }
           || !type.IsDefined(typeof(TenancyDbFunctionsAttribute), inherit: false);
}

[TenancyDbFunctions]
public static class TenancyDbFunctionsExtensions { /* CurrentTenant() */ }
\end{sourcecode}

The plugin lives in the \textbf{core package} and recognizes function classes by an
attribute. A single plugin thus also covers the PostgreSQL-specific functions from
chapters~\ref{ch:custom-sql-expression-nodes} and \ref{ch:aggregate-translators} and every
future provider package, without the core having to know these classes.

\begin{important}[When the plugin is needed]
  Whether a custom EF.Functions method needs this plugin depends on whether it has arguments
  referencing data.
  \code{EF.\allowbreak{}Functions.\allowbreak{}Soundex(\allowbreak{}c.\allowbreak{}Name)} is
  harmless because \code{c.\allowbreak{}Name} points to a column. Methods without arguments
  or only with constants (\code{Current\allowbreak{}Tenant(\allowbreak{})},
  \code{Random(\allowbreak{})}, \code{Server\allowbreak{}Time(\allowbreak{})}) always need
  it. Without it, the exception comes from the method body, which is hard to diagnose.
\end{important}

% ------------------------------------------------------------
\chapter{Further query extension points}\label{ch:further-query-extension-points}

The sample needs method and member translators. Complex plugins use further points; here is
an overview with classification:

\begin{datatable}{H{0.96}H{1.08}H{0.96}}
  \textbf{Extension point} & \textbf{Purpose} & \textbf{Notes} \\
  \tablehead
  \code{IAggregate\allowbreak{}Method\allowbreak{}Call\allowbreak{}Translator\allowbreak{}Plugin}
  & Custom aggregate functions in \code{Group\allowbreak{}By} (e.g.
  \code{bool\_\allowbreak{}and}, statistics functions) & Scoped, multiple registrations;
  implemented in chapter~\ref{ch:aggregate-translators} \\
  \tablerule
  \code{IQuery\allowbreak{}Expression\allowbreak{}Interceptor} & Rewrite the LINQ expression
  tree before translation & Singleton interceptor; ideal for query rewrites, tags, hints \\
  \tablerule
  \code{IMaterialization\allowbreak{}Interceptor} & Create, initialize or check objects
  during materialization & Singleton interceptor; must be stateless and thread-safe \\
  \tablerule
  \code{IMethod\allowbreak{}Call\allowbreak{}Translator\allowbreak{}Provider} etc. & Replace
  translator infrastructure & Single provider services; only in exceptional cases \\
  \tablerule
  \code{IQuery\allowbreak{}Sql\allowbreak{}Generator\allowbreak{}Factory} & SQL generation
  for custom SqlExpression types & Singleton, provider service; implemented in
  chapter~\ref{ch:custom-sql-expression-nodes} \\
  \tablerule
  \code{IRelational\allowbreak{}Parameter\allowbreak{}Based\allowbreak{}Sql\allowbreak{}Processor\allowbreak{}Factory}
  & Nullability processing of custom expression nodes & Scoped, provider service;
  implemented in chapter~\ref{ch:custom-sql-expression-nodes} \\
  \tablerule
  \code{IQuery\allowbreak{}Translation\allowbreak{}Postprocessor\allowbreak{}Factory} &
  Adjust the expression tree after translation & Provider service; deep intervention,
  tightly bound to EF internals \\
\end{datatable}

\begin{note}[Precompiled queries]
  With NativeAOT and precompiled queries (since EF Core 9), EF translates queries at build
  time. Translators then run in the build process, not at run time. They may therefore only
  depend on options and model, never on runtime state. The sample meets this: the setting
  name is an option; the tenant itself is only evaluated in the database session.
\end{note}
